\subsection{Empirical evidence argues against symmetric structural equivalence}

The directional witness is not only a conservative formal choice. Wang reports a controlled asymmetry in structural transfer between natural-language and biological foundation models \cite{wang2026asymmetry}. Language models trained or fine-tuned on structural language tasks retain meaningful competence when moved to biological-sequence structure, while the reverse transfer remains weak across matched-token controls, scaling and multiple biological model families. The result is domain-specific and does not establish a universal law of transfer, but it is a direct counterexample to the assumption that shared structural regularity implies symmetric representational transfer.

For Paper 3 this has two consequences. First, a structural relation should default to a directed edge $S_A\to S_B$ with a witness rather than an undirected equivalence class. Second, transfer evaluation should include direction reversal as an explicit perturbation. If $A\to B$ succeeds but $B\to A$ fails, the library should preserve that asymmetry rather than canonicalize the two tasks into one globally interchangeable class. Any later transitive composition likewise requires a new scope/invariant check.

This also strengthens the Q2/Q3 benchmark logic. The object being tested is not whether two domains ``look analogous'' in a human narrative sense. It is whether a particular source-side operator or inference depends only on relations and invariants that are actually available in the target under the registered boundary conditions. Directional negative-transfer history is therefore part of the reusable state.
